% =============================================================================
% Capítulo 6 · Conclusiones preliminares y recomendaciones (iteración 3c, 7-nov-2026)
% Cada conclusión se ata a un objetivo o hipótesis (cap. 1) y a la evidencia del cap. 5 o a un
% aporte formal del cap. 3. Sin cifras nuevas: las que se mencionan entran por macro.
% =============================================================================
\chapter{Conclusiones preliminares y recomendaciones}
\label{cap:conclusiones}

Este capítulo reúne lo que la tesis puede afirmar en su estado actual. Las conclusiones son
\textbf{preliminares}: se apoyan en el avance de resultados del \cref{cap:resultados}, que aún no se
ha discutido con el asesor, y en los aportes formales del \cref{cap:marco}, que sí están
demostrados. Cada una indica el objetivo o la hipótesis a la que responde y la evidencia que la
sostiene. La \cref{sec:conc-recomendaciones} formula las recomendaciones y los trabajos futuros, y
la \cref{tab:estado-objetivos} resume el estado de cada objetivo.

\section{Conclusiones preliminares}
\label{sec:conc-conclusiones}

\begin{enumerate}[label=\textbf{C\arabic*.}, ref=C\arabic*, leftmargin=*]
  \item\label{conc:datos} \textbf{Datos trazables y reproducibles} (\cref{oe:datos}). La tesis
    dispone de \nSeries{} series mensuales del BCRP, de \fechaInicioNiveles{} a \fechaFinDatos{},
    congeladas e identificadas por su huella (decisión D4), con exclusiones justificadas por la
    estructura de los datos y no por los resultados (decisiones D1 y D16). Todo el análisis se
    calcula en Python y en R con paridad verificada (\cref{sec:met-implementacion,ap:reproducibilidad}).
  \item\label{conc:formal} \textbf{La cadena de la escuela de Camiz queda formalizada}
    (\cref{oe:clasicos,oe:ventana}). El \cref{cap:marco} da definiciones y supuestos explícitos a
    la HFC, al EPCA con la función \texttt{aede}, a la partición de Fisher con longitud mínima y a
    los dos nulos por remuestreo, y demuestra las propiedades que la tesis usa: que la correlación
    ponderada de una ventana es una matriz de correlaciones solo si la estandarización también es
    ponderada (\cref{prop:epca,obs:m1}); que el tamaño efectivo de los pesos gaussianos es una
    fracción del ancho que no depende de los datos (\cref{prop:kish}); que el RV de cargas compara
    aproximaciones de rango $d$ y no depende de los signos ni del orden de los ejes
    (\cref{prop:rv-cargas}); que la longitud mínima puede romper la monotonía del costo de Fisher
    (\cref{prop:fisher-monotonia}); que la variante del índice de Caliński y Harabasz de
    \texttt{aede} no elige menos segmentos que la habitual (\cref{prop:ch}), y que el criterio
    de ancho de ventana es válido bajo intercambiabilidad (\cref{prop:oe4}). Estos resultados son
    el aporte formal de la tesis, con independencia de lo que muestren los datos.
  \item\label{conc:niveles} \textbf{Hay razones formales para no diferenciar}
    (\cref{sec:epistemologia}). La primera diferencia atenúa las oscilaciones lentas que definen el
    ciclo financiero (\cref{prop:ganancia}), y en los datos la proporción de varianza en
    oscilaciones de dos años o más cae fuertemente al transformar las series
    (\cref{sec:ap-no-diferenciar}). Analizar las series sin transformar es coherente con el objeto
    de estudio, pero obliga a separar la tendencia de la estructura (\ref{conc:tendencia}).
  \item\label{conc:pca} \textbf{Tres ejes interpretables y estables} (\cref{oe:clasicos,oe:dimension}).
    El PCA de las series sin transformar tiene tres ejes con lectura económica clara:
    profundización, tasas de interés de corto plazo y riesgo de mercado; los tres son
    descriptivamente estables bajo remuestreo y el cuarto no lo es (\cref{sec:res-pca}). Los
    criterios de retención no coinciden en el tercer eje, de modo que $d=3$, fijado de antemano
    (decisión D9), se sostiene con esa salvedad.
  \item\label{conc:tendencia} \textbf{El primer eje es la tendencia común; las tasas, un
    comovimiento genuino} (\cref{oe:clasicos}). El primer eje es casi exactamente una tendencia
    temporal, y su inercia no supera la de series independientes con su propia tendencia (nulo
    mixto, decisiones D8 y D19). El eje de tasas sí la supera (\cref{sec:res-tendencia}). En
    niveles, por tanto, el primer componente no debe leerse como un índice del ciclo financiero,
    sino como profundización y crecimiento nominal.
  \item\label{conc:hfc} \textbf{Bloques estables que reproducen los del antecedente} (HT2;
    \cref{oe:clasicos,oe:robustez-hfc}). La HFC separa \clasesHfcMenor{} bloques (profundización,
    tasas y riesgo de mercado) que se reproducen casi siempre bajo remuestreo por bloques y se
    corresponden, por bloque, con los de \textcite{camiz_cerda_atoche_2026}
    (\cref{sec:res-hfc}). HT2 queda sostenida en lo esencial.
  \item\label{conc:oe4} \textbf{El criterio de ventana está formalizado, pero en niveles no
    discrimina} (HT1; \cref{oe:ventana}). El criterio de OE4 se definió, se demostró su validez en
    una versión simétrica y se calculó para los cuatro anchos. Con las series sin transformar,
    ningún ancho muestra más cambios de estructura de los que produce una estructura constante
    (\cref{sec:res-oe4}), y la transición de octubre de 2017 no se confirma
    (\cref{sec:res-transiciones}). HT1 no queda sostenida por el criterio formal en niveles; la
    causa más probable es la escasa potencia del nulo con series integradas, que debe evaluarse.
  \item\label{conc:datacion} \textbf{Regímenes alrededor de los grandes choques, con un número
    que depende de la regla} (HT3; \cref{oe:datacion,oe:episodios}). La datación principal separa
    un tramo alrededor de la crisis financiera global, un tramo intermedio largo y un tramo
    posterior a la pandemia, y los dos choques más intensos se reflejan en ella y, la pandemia,
    también en el EPCA. Pero el número de regímenes cambia con el máximo de la grilla, uno de los
    cortes es poco estable y los episodios intermedios no se separan (\cref{sec:res-datacion,sec:res-validacion}).
    HT3 queda parcialmente sostenida.
  \item\label{conc:dl} \textbf{El complemento de aprendizaje profundo está formalizado, pero sin
    evaluar} (HT4; \cref{oe:dl,oe:incertidumbre,oe:comparacion}). La \cref{prop:ae-lineal} ancla el
    autoencoder en el PCA, y el diseño de anclaje lineal, validación por bloques y estabilidad entre
    réplicas está definido (\cref{sec:dl}); pero los modelos aún no se han entrenado con las
    \nSeries{} series vigentes (decisión D17.4). HT4 queda sin evaluar.
\end{enumerate}

En síntesis, con las series sin transformar la tesis puede afirmar qué estructura tienen los datos
(una tendencia común, un comovimiento de tasas y bloques estables) y dar a la cadena de métodos un
fundamento formal que no tenía; todavía no puede afirmar que los cambios de esa estructura en el
tiempo sean distinguibles de los que produciría una estructura constante.

\section{Recomendaciones y trabajos futuros}
\label{sec:conc-recomendaciones}

\paragraph{Recomendaciones metodológicas (a discutir con el asesor).}

\begin{enumerate}[label=\textbf{R\arabic*.}, ref=R\arabic*, leftmargin=*]
  \item Discutir la Parte~C de la decisión D2 con los resultados de este capítulo a la vista:
    en particular, si la pregunta de HT1 y el criterio de OE4 deben evaluarse en una vía en que la
    variabilidad de la correlación entre ventanas sea menor (\ref{conc:oe4}).
  \item Separar la tendencia de la estructura en la datación: evaluar, como sensibilidad, la
    partición de trayectorias que no incluyan el primer eje o que usen los \emph{scores} sin
    tendencia (\ref{conc:tendencia}, \ref{conc:datacion}). Es una propuesta, no una decisión
    tomada.
  \item Medir la potencia del criterio de OE4 con un contexto simulado de estructura conocida
    (\cref{oe:simulado}), con series integradas como las del BCRP: si el criterio no detecta
    cambios simulados de magnitud comparable, su resultado en niveles no informa sobre HT1.
  \item Completar el contraste externo con una tabla de compatibilidad entre cortes y episodios y
    con episodios posteriores a la pandemia fechados con fuentes externas al análisis (decisión
    D18).
\end{enumerate}

\paragraph{Trabajos futuros.} Volver a entrenar los autoencoders y las redes LSTM con las
\nSeries{} series (iteración 4); explorar los objetivos condicionales (la HFC evolutiva, el STATIS
dual como detector independiente y la extensión de la muestra a años anteriores a
\anioInicioDatos{}), y completar la revisión de la literatura peruana. El \cref{cap:pendientes}
detalla cada tema, su causa y de quién depende.

\paragraph{Recomendación de uso.} Un indicador de régimen transparente y reproducible podría
servir como insumo exploratorio de los análisis de estabilidad financiera (\cref{sec:justificacion}).
En su estado actual, el resultado no debe usarse con ese fin: el primer eje en niveles mide
profundización y no ciclo, y los cambios de estructura no se han distinguido todavía de la
variabilidad muestral.

\begin{table}[htbp]
\centering
\small
\caption{Estado de los objetivos específicos en esta versión}
\label{tab:estado-objetivos}
\begin{tabular}{@{}l>{\raggedright\arraybackslash}p{0.5\textwidth}>{\raggedright\arraybackslash}p{0.3\textwidth}@{}}
\toprule
Objetivo & Estado & Dónde \\
\midrule
\ref{oe:estado-arte} & Avanzado; literatura peruana incompleta & \cref{cap:estado-arte} \\
\ref{oe:datos} & Cumplido & \cref{sec:met-datos}, \cref{ap:reproducibilidad} \\
\ref{oe:clasicos} & Avanzado en niveles; Parte~C pendiente & \cref{sec:res-pca,sec:res-hfc,sec:res-dinamica} \\
\ref{oe:ventana} & Formalizado y calculado; no discrimina en niveles & \cref{sec:criterio-oe4}, \cref{sec:res-oe4} \\
\ref{oe:robustez-hfc} & Cumplido & \cref{sec:res-hfc} \\
\ref{oe:dimension} & Parcial: falta formalizar el criterio de Rencher & \cref{sec:retencion}, \cref{sec:res-pca} \\
\ref{oe:dl} & Diseñado; sin entrenar con \nSeries{} series & \cref{sec:dl}, \cref{sec:res-dl} \\
\ref{oe:incertidumbre} & Pendiente & \cref{sec:pend-dl} \\
\ref{oe:datacion} & Avanzado con el PCA; fases preliminares & \cref{sec:res-datacion} \\
\ref{oe:episodios} & Avanzado; falta la tabla de compatibilidad & \cref{sec:res-validacion} \\
\ref{oe:simulado} & No iniciado & \cref{sec:pend-metodo} \\
\ref{oe:comparacion} & Pendiente de \ref{oe:dl} & \cref{sec:pend-dl} \\
\bottomrule
\end{tabular}
\end{table}
